%%
%% Berliner Hochschule für Technik --  Abschlussarbeit
%%
%% Kapitel 6
%%
%%

\chapter{Planung und Durchführung der App-Studie}

\section{Zielsetzung und Studienlogik}

Die App-Studie dient der Prüfung, ob der im Labor beschriebene kurzfristige Unterschied zwischen Cue-Matching und Cue-Labeling auch unter Alltagsbedingungen mit einer mobilen Android-App beobachtet werden kann. Untersucht wird nicht ein langfristiger Rauchstopp, sondern das unmittelbar berichtete subjektive Rauchverlangen nach der Bearbeitung rauchbezogener Auslöserreize. Die App ist deshalb als Studienprototyp und nicht als Rauchentwöhnungsprogramm oder medizinische Behandlung zu verstehen. Diese Abgrenzung ist methodisch und ethisch wesentlich, weil die Studie zwar mit gesundheitsbezogenen Angaben arbeitet, aber keinen therapeutischen Anspruch erhebt.

Die zentrale abhängige Variable ist das selbstberichtete Rauchverlangen auf einer Skala von 0 bis 100. Die Skala ist für die mobile Nutzung geeignet, weil sie ohne ausführliches Instrumentarium verstanden werden kann, auf kleinen Displays gut abbildbar ist und numerisch unmittelbar ausgewertet werden kann. Gleichzeitig erlaubt sie eine feinere Abstufung als die im zugrunde liegenden Laborversuch verwendete fünfstufige Skala. Für den Vergleich mit der Laborarbeit kann eine lineare Umrechnung erfolgen. Der dort berichtete Mittelwertunterschied von 0,1 Punkten auf einer Skala von 1 bis 5 entspricht auf der hier verwendeten Skala einem Unterschied von
\begin{equation}
    \Delta = 0{,}1 \times 25 = 2{,}5
\end{equation}
Punkten. Die ebenfalls umgerechnete Standardabweichung der paarweisen Differenzen beträgt
\begin{equation}
    s_D = 0{,}3279 \times 25 = 8{,}1968.
\end{equation}
Damit wird die Studie bewusst auf einen kleinen, kurzfristigen Effekt ausgerichtet. Eine mögliche statistische Signifikanz wäre daher nicht automatisch mit einem klinisch bedeutsamen Nutzen gleichzusetzen. Für die Interpretation wird zwischen einer wahrnehmbaren Veränderung und einem für Betroffene tatsächlich wertvollen Nutzen unterschieden \cite{tabibnia_cue_labeling_2026,weinfurt_clarifying_2019}.

\section{Studiendesign}

Geplant ist eine quantitative Pilotstudie mit Primärdatenerhebung und Within-Subject-Design. Jede teilnehmende Person bearbeitet sowohl Cue-Matching- als auch Cue-Labeling-Aufgaben. Dadurch werden interindividuelle Unterschiede im Rauchverlangen, in der Nikotinabhängigkeit, im Tagesrhythmus und in der Reaktion auf Bildreize reduziert, weil jede Person als eigene Vergleichsbasis dient.

Die App führt die Teilnehmenden über einen Zeitraum von bis zu zwei Wochen durch 20 Studiensituationen. In jeder Studiensituation werden rauchbezogene Bilder gezeigt. Anschließend ordnen die Teilnehmenden entweder einem Auslöserreiz ein anderes rauchbezogenes Bild zu oder sie wählen eine passende textliche Beschreibung des gezeigten Auslöserreizes. Nach der Aufgabenbearbeitung wird das aktuelle Rauchverlangen auf der Skala von 0 bis 100 abgefragt. Die geplante Aufgabenverteilung sieht insgesamt 50 Zuordnungen im Cue-Matching, 50 Zuordnungen im Cue-Labeling und 20 Selbstberichte zum Rauchverlangen vor.

Zwischen zwei Studiensituationen liegt ein technischer Mindestabstand von mehreren Stunden. Dadurch sollen unmittelbar aufeinanderfolgende Selbstberichte vermieden werden. Die Vorgabe ist zugleich alltagstauglich: Die Teilnehmenden müssen nicht an festen Laborterminen erscheinen und können die Nutzung in ihren Tagesablauf integrieren. Obwohl die Studie für zwei Wochen vorgesehen ist, kann die vollständige Bearbeitung bei regelmäßiger Nutzung bereits nach ungefähr sieben Tagen abgeschlossen sein. Optionale Benachrichtigungen erinnern an die Nutzung, dürfen die Teilnehmenden aber nicht zu einer unmittelbaren Bearbeitung drängen.

Die Reihenfolge der Bedingungen wird innerhalb der App so gesteuert, dass beide Aufgabentypen in der Stichprobe möglichst gleichmäßig auftreten. Praktisch kann dies durch eine vorbereitete Liste von 20 Studiensituationen mit zehn Cue-Matching- und zehn Cue-Labeling-Blöcken erfolgen, deren Reihenfolge pro Installation zufällig permutiert wird. Diese Vorgehensweise ist technisch robust, weil sie ohne serverseitige Sitzungssteuerung auskommt, und methodisch ausreichend, weil die Bedingungsverteilung pro teilnehmender Person ausgeglichen bleibt. Die App speichert lokal, welche Studiensituationen bereits abgeschlossen wurden, und übermittelt nur die für die Studienauswertung erforderlichen Selbstberichte.

\section{Einschluss, Rekrutierung und Teilnahmeprüfung}

Eingeschlossen werden Personen, die zwischen 30 und 65 Jahre alt sind, derzeit mindestens zehn Zigaretten pro Tag rauchen, Deutsch oder Englisch lesen können, dauerhaft in Deutschland oder in der Schweiz leben und ein Android-Smartphone oder Android-Tablet mit Internetzugang besitzen. Zusätzlich müssen sie bereit sein, die App im Studienzeitraum zu verwenden und insgesamt 20 Selbstberichte abzugeben. Jede Person darf nur einmal teilnehmen.

Die Teilnahmebedingungen werden anhand der Selbstauskunft geprüft. Eine Ausweisprüfung, eine Kohlenmonoxidmessung der Atemluft oder ein Cotinintest werden nicht durchgeführt. Diese Entscheidung reduziert Eingriffsintensität, organisatorischen Aufwand und die Verarbeitung zusätzlicher sensibler Daten. Sie senkt allerdings die Sicherheit der Prüfung. Deshalb werden die Selbstauskünfte und Selbstberichte nachträglich statistisch plausibilisiert. Auffällige Muster können dokumentiert und bei der Interpretation berücksichtigt werden. Die einmalige Teilnahme wird praktikabel über organisatorische Daten wie Name, E-Mail-Adresse und Bankverbindung sowie über die Anzahl der pro App-Installation übermittelten Selbstberichte kontrolliert.

Die Rekrutierung kann über Online-Plattformen, thematisch passende Online-Communities sowie Flyer oder Aushänge in Apotheken, Arztpraxen und Beratungsstellen erfolgen. In allen Rekrutierungsmaterialien muss deutlich erkennbar sein, dass es sich um eine wissenschaftliche Studie mit einem App-Prototyp handelt. Formulierungen, die eine therapeutische Wirksamkeit, eine Unterstützung beim Rauchstopp oder einen individuellen Gesundheitsnutzen versprechen, werden vermieden. Zulässig und angemessen ist die Darstellung, dass untersucht wird, ob sich subjektives Rauchverlangen nach rauchbezogenen Bildern kurzfristig unterschiedlich entwickelt, wenn Reize beschriftet oder mit anderen Reizen verglichen werden.

\section{Ablauf der Teilnahme}

Vor der Nutzung erhalten Interessierte eine Studieninformation, eine Datenschutzerklärung und eine Einwilligungserklärung. Die Studieninformation beschreibt Ziel, Verantwortlichkeit, Teilnahmebedingungen, Ablauf, Datenkategorien, lokale Bildverarbeitung, mögliche Belastungen, Freiwilligkeit, Abbruchmöglichkeit und Aufwandsentschädigung. Erst nach Einwilligung und Prüfung der Teilnahmebedingungen wird die App bereitgestellt beziehungsweise zur Nutzung freigeschaltet.

Für Installation und Nutzung ist eine Internetverbindung erforderlich. Die Größe der eigentlichen Studienübermittlungen ist gering, weil im Studienbetrieb nur Selbstberichte und technische Zählerstände übertragen werden. Die App-Installation kann abhängig vom gewählten Funktionsumfang größer sein, insbesondere wenn KI-Modelle auf dem Gerät bereitgestellt werden. Deshalb wird den Teilnehmenden empfohlen, Installation und Nutzung nach Möglichkeit über WLAN vorzunehmen.

Eine Studiensituation besteht aus vier Schritten. Erstens prüft die App, ob die nächste Studiensituation freigegeben ist und der Mindestabstand zur vorherigen Situation eingehalten wurde. Zweitens präsentiert sie die vorgesehenen rauchbezogenen Bilder. Drittens bearbeiten die Teilnehmenden je nach Bedingung eine Matching- oder Labeling-Aufgabe. Viertens geben sie ihr aktuelles Rauchverlangen auf der Skala von 0 bis 100 an. Nach erfolgreicher Übermittlung wird der lokale Fortschritt aktualisiert. Sind 20 verwertbare Selbstberichte abgegeben, ist die Teilnahme vollständig.

Die Teilnahme ist freiwillig und kann jederzeit abgebrochen werden. Ein Abbruch führt zu keinen Nachteilen. Da die Aufwandsentschädigung an eine vollständige und verwertbare Teilnahme gebunden ist, wird sie nur ausgezahlt, wenn die Teilnahmebedingungen erfüllt sind und alle erforderlichen Selbstberichte vorliegen. Vorgesehen sind 15,00 Euro für Teilnehmende aus Deutschland und 15,00 Schweizer Franken für Teilnehmende aus der Schweiz. Bei technischen Problemen sollen Teilnehmende zeitnah Kontakt aufnehmen; falls die Ursache nachvollziehbar technisch ist, kann ein angemessener Aufschub für die Übermittlung der Berichte gewährt werden.

\section{Datenflüsse und Datenschutz im Studienbetrieb}

Der Studienbetrieb folgt dem Grundsatz der Datensparsamkeit. Für Organisation und Abrechnung werden Name, E-Mail-Adresse und Bankverbindung benötigt. Für die Prüfung der Teilnahmebedingungen werden Alter, Wohnsitzland, Sprachverständnis, Android-Endgerät und Rauchstatus verarbeitet. Für die wissenschaftliche Auswertung werden die Selbstberichte zum Rauchverlangen, die Bedingung der jeweiligen Studiensituation und technisch notwendige Fortschrittsinformationen verwendet. Die wissenschaftliche Auswertung erfolgt auf Grundlage anonymisierter Studienergebnisse; Kontakt- und Zahlungsdaten werden getrennt von den anonymisierten Ergebnisdaten behandelt.

Die App und der Studienserver tauschen im Normalbetrieb nur Angaben zum Rauchverlangen und zur Anzahl bisher übermittelter Selbstberichte aus. Eigene Bilddateien oder Kamerabilder können, sofern diese Funktion aktiviert ist, lokal auf dem Endgerät für eine Bilderkennung verwendet werden. Diese Bilder werden nicht an den Studienserver übertragen, nicht dauerhaft gespeichert und nach der jeweiligen Verarbeitung verworfen. Diese Architektur reduziert das Risiko für die Teilnehmenden erheblich, weil potentiell besonders sensible Alltagsbilder das Endgerät nicht verlassen.

Aus informatischer Sicht ist die Trennung der Datenflüsse ein zentrales Entwurfsprinzip. Ergebnisdaten werden so strukturiert, dass sie für die statistische Auswertung ausreichen, aber keine unmittelbare Identifikation der Teilnehmenden ermöglichen. Organisatorische Daten werden dagegen nur für Kontakt, Teilnahmeprüfung und Auszahlung verwendet. Die App begrenzt die Anzahl der übermittelbaren Selbstberichte technisch auf 20 und verhindert dadurch Mehrfachübermittlungen derselben Installation. Der Server muss Eingaben validieren, Wertebereiche prüfen und unvollständige oder widersprüchliche Datensätze kennzeichnen, ohne dadurch unnötige Zusatzdaten zu erzeugen.

\section{Studienbetrieb, Monitoring und Support}

Während der Durchführung wird die technische Infrastruktur überwacht. Dazu gehören die Erreichbarkeit des Servers, die erfolgreiche Annahme von Selbstberichten, Plausibilitätsprüfungen der Wertebereiche und die Erkennung auffälliger Übermittlungsmuster. Das Monitoring dient nicht der Verhaltenskontrolle einzelner Teilnehmender, sondern der Sicherstellung, dass die Studie technisch ordnungsgemäß durchgeführt werden kann.

Supportanfragen werden vor allem bei Installation, fehlgeschlagenen Übermittlungen, Problemen mit Benachrichtigungen oder Missverständnissen zum Studienablauf erwartet. Der Support muss so gestaltet werden, dass er die Studiendaten nicht verfälscht. Inhaltliche Hinweise, die Teilnehmende zu bestimmten Antworten auf der Rauchverlangensskala bewegen könnten, sind zu vermeiden. Zulässig sind technische Hilfen, Erläuterungen zum Ablauf und Hinweise auf die Möglichkeit, die Teilnahme zu unterbrechen oder abzubrechen.

Ein besonderes Risiko besteht darin, dass rauchbezogene Bilder kurzfristig Rauchverlangen auslösen oder verstärken. Deshalb informiert die Studie vorab über diese Belastung. Teilnehmende werden darauf hingewiesen, die Nutzung zu unterbrechen oder abzubrechen, wenn sie die App als unangenehm erleben. Bei starkem Suchtdruck, gesundheitlichen Beschwerden oder psychischer Belastung wird auf medizinische, psychologische oder suchtberatende Unterstützung verwiesen. Diese Hinweise sind Bestandteil einer verantwortbaren Studiendurchführung, weil der Prototyp keine Behandlung ersetzt.

\section{Qualitätssicherung und Auswertungsvorbereitung}

Die Auswertung setzt voraus, dass die 20 Selbstberichte pro teilnehmender Person vollständig, eindeutig einer Bedingung zugeordnet und im zulässigen Wertebereich von 0 bis 100 liegen. Bereits bei der Datenerfassung werden deshalb technische Validierungen eingesetzt. Werte außerhalb des zulässigen Bereichs, fehlende Bedingungskennzeichnungen, doppelte Übermittlungen oder offensichtlich unvollständige Datensätze werden serverseitig zurückgewiesen oder markiert.

Für die spätere Analyse werden pro Person paarweise Differenzen zwischen Cue-Labeling- und Cue-Matching-Situationen gebildet. Der geplante Haupteffekt ist die Differenz des berichteten Rauchverlangens nach beiden Aufgabenarten. Für die Stichprobenplanung wird der aus der Laborarbeit abgeleitete standardisierte Effekt verwendet:
\begin{equation}
    d = \frac{2{,}5}{8{,}1968} \approx 0{,}305.
\end{equation}
Bei einem gepaarten Test, einer einseitigen Alternativhypothese und einer angestrebten Teststärke von 80\,\% ergibt sich daraus eine Zielgröße von 68 verwertbaren Teilnehmenden. Da bei einer App-Studie mit Abbrüchen, technischen Problemen und unvollständigen Datensätzen zu rechnen ist, wird ein zusätzlicher Ausfallpuffer von etwa 25\,\% vorgesehen. Daraus ergibt sich ein Rekrutierungsziel von ungefähr 85 Teilnehmenden \cite{cohen_statistical_2013,giner_sorolla_power_2024}.

Die Qualitätssicherung umfasst darüber hinaus die Dokumentation der App-Version, des Zeitpunkts der Datenerhebung, der verwendeten Aufgabenlogik und der Datenbereinigung. Diese Angaben sind notwendig, um die Ergebnisse nachvollziehbar zu interpretieren. Insbesondere muss erkennbar bleiben, ob unvollständige Datensätze wegen technischer Fehler, Abbruch der Teilnahme oder Nichterfüllung der Teilnahmebedingungen ausgeschlossen wurden.

\section{Methodische und ethische Einordnung}

Die Studie verbindet eine medizinische Fragestellung mit einer informatischen Umsetzung. Der medizinische Beitrag liegt in der alltagsnahen Prüfung eines kurzfristigen verhaltensmedizinischen Effekts auf Rauchverlangen. Der informatische Beitrag liegt in der Gestaltung eines datensparsamen mobilen Studiensystems, das Aufgabensteuerung, lokale Verarbeitung, Fortschrittskontrolle, sichere Übermittlung und auswertbare Datenstrukturen miteinander verbindet.

Die Aussagekraft bleibt durch den Pilotcharakter begrenzt. Die Teilnahmebedingungen werden nicht klinisch überprüft, die Nutzung findet auf unterschiedlichen Android-Endgeräten statt, und Alltagskontexte können nur begrenzt kontrolliert werden. Diese Einschränkungen sind jedoch nicht nur methodische Schwächen, sondern Teil der Forschungsfrage: Die App soll gerade unter realen Bedingungen funktionieren. Deshalb werden Robustheit, Datenschutz, geringe Belastung und nachvollziehbare Auswertung höher gewichtet als maximale Kontrolle.

Aus ethischer Sicht ist entscheidend, dass die Teilnehmenden vorab über Zweck, Ablauf, Datenverarbeitung, Risiken und Grenzen der App informiert werden. Die App darf Rauchverlangen nicht verharmlosen und keine medizinische Wirkung versprechen. Ihre Rolle ist die eines Studieninstruments zur Erhebung und experimentellen Variation kurzfristiger Reaktionen auf Auslöserreize. Erst die spätere Auswertung kann zeigen, ob Cue-Labeling im Alltag einen messbaren Unterschied gegenüber Cue-Matching erzeugt und ob dieser Unterschied über eine statistische Beobachtung hinaus praktisch bedeutsam sein könnte.
